\chapter{परिचक्रीय अभिक्रियाएँ}\label{ch:b3:pericyclic}

त्वचा पर पड़ता सूर्यप्रकाश एक सेकंड के अंश में 7-डीहाइड्रोकोलेस्टेरॉल का एक वलय खोल देता है;
फिर शरीर की ऊष्मा एक हाइड्रोजन परमाणु को सात कार्बन परमाणुओं के पार खिसका देती है,
और उत्पाद विटामिन D$_3$ है। किसी भी चरण में मध्यवर्ती नहीं होता: आबंध परमाणुओं के एक वलय के
चारों ओर, एक ही संक्रमण अवस्था से होकर, साथ-साथ टूटते और बनते हैं। पहला
चरण केवल प्रकाश से होता है और दूसरा केवल ऊष्मा से, और प्रत्येक
एक ही त्रिविम समावयवी देता है। यह अध्याय संबंधित आण्विक कक्षकों की सममिति से
इसका कारण समझाता है: रॉबर्ट वुडवर्ड और रोआल्ड
हॉफ़मैन के नियम, जो पूर्वानुमान करते हैं कि ऐसी अभिक्रिया हो सकती है या नहीं और वह क्या देती है।

\begin{recall}
द्वितीय वर्ष के खंड ने पॉलिईनों के ह्युकेल गुणांक, ऐरोमैटिक और
प्रति-ऐरोमैटिक वलय, HOMO, LUMO और अग्रांत कक्षक, समन्वित अभिक्रियाएँ, और
डाईनरागी के साथ डाईन की डील्स--ऐल्डर अभिक्रिया, उसके endo नियम सहित, दिए; प्रथम
वर्ष के खंड ने \textit{Z}/\textit{E} वर्णक। \Cref{ch:b3:point-groups} ने $C_2$ अक्ष
और \omterm{def:b3:point-groups:mirror}{दर्पण तल} परिभाषित किए, \cref{ch:b3:photochemistry} ने उत्तेजित अवस्थाएँ और [2+2]
प्रकाश-चक्रीयोजन।
\end{recall}

\begin{omfigure}
\includegraphics[width=0.55\linewidth]{images/book4/ai/b3-26-terrace.jpg}

{\small धूप भरी छत पर पढ़ते हुए: त्वचा तक पहुँचने वाला पराबैंगनी प्रकाश एक
परिचक्रीय वलय-विवरण द्वारा विटामिन D$_3$ का संश्लेषण आरंभ करता है।}
\end{omfigure}

\section{चार परिवार}

\begin{definition}[परिचक्रीय अभिक्रिया]\label{def:b3:pericyclic:pericyclic}
\emph{परिचक्रीय अभिक्रिया}\index{परिचक्रीय अभिक्रिया} ऐसी समन्वित अभिक्रिया है
जिसकी संक्रमण अवस्था अन्योन्यक्रिया करते कक्षकों का एक चक्रीय विन्यास होती है, जिसमें सभी
आबंध वलय के चारों ओर एक ही समय पर बनते और टूटते हैं।
\end{definition}

\begin{definition}[परिचक्रीय अभिक्रियाओं के परिवार]\label{def:b3:pericyclic:families}
\emph{विद्युतचक्रीय अभिक्रिया}\index{विद्युतचक्रीय अभिक्रिया} में एक संयुग्मित
पॉलिईन अपने सिरों के बीच एक $\sigma$ आबंध बनाकर वलय में बंद होता है, या
इसका उलटा। \emph{चक्रीयोजन}\index{चक्रीयोजन} में दो $\pi$ निकाय
दो $\sigma$ आबंध बनाकर एक वलय में जुड़ते हैं; $[m + n]$ चक्रीयोजन
$m$ और $n$ परमाणुओं (या इलेक्ट्रॉनों) के घटकों को जोड़ता है। \emph{सिग्माट्रॉपी
पुनर्विन्यास}\index{सिग्माट्रॉपी पुनर्विन्यास} में एक $\sigma$ आबंध किसी $\pi$
निकाय के अनुदिश प्रवास करता है; $[i, j]$ विस्थापन में आबंध दोनों अणुखंडों के परमाणुओं 1 और 1 से परमाणुओं $i$ और $j$
पर जाता है। \emph{कीलेट्रॉपी अभिक्रिया}\index{कीलेट्रॉपी अभिक्रिया}
एक ही परमाणु से दो आबंध बनाती या तोड़ती है, जैसे जब सल्फ़र डाइऑक्साइड
किसी सल्फ़ोलीन से निकलती है। \emph{ईन अभिक्रिया}\index{ईन अभिक्रिया} में
ऐलिलिक हाइड्रोजन वाला एक ऐल्कीन किसी $\pi$ आबंध से जुड़ता है, और हाइड्रोजन उसके साथ खिसकता है।
\end{definition}

\begin{definition}[समफलकीय और विपरीतफलकीय]\label{def:b3:pericyclic:faciality}
कोई घटक \emph{समफलकीय}\index{समफलकीय} रूप से अभिक्रिया करता है जब उसके दो नए
आबंध बनते हैं, या दो पुराने टूटते हैं, उसके $\pi$ निकाय के एक ही फलक पर (या
$\sigma$ आबंध पर धारण के साथ), और \emph{विपरीतफलकीय}\index{विपरीतफलकीय} रूप से जब वे
विपरीत फलकों पर हों (या प्रतिलोमन के साथ)।
\end{definition}

\section{विद्युतचक्रीय अभिक्रियाएँ}

\begin{definition}[सहघूर्णी और विघूर्णी]\label{def:b3:pericyclic:rotation-modes}
विद्युतचक्रीय वलय समापन में दोनों अंतस्थ समूह उन आबंधों के परितः घूमते हैं
जो उन्हें शृंखला से जोड़ते हैं। घूर्णन \emph{सहघूर्णी}\index{सहघूर्णी} होता है
जब दोनों एक ही दिशा में घूमें (शृंखला के अनुदिश देखने पर दोनों दक्षिणावर्त), और
\emph{विघूर्णी}\index{विघूर्णी} जब वे विपरीत दिशाओं में घूमें।
\end{definition}

\begin{theorem}[विद्युतचक्रीय अभिक्रियाएँ]\label{thm:b3:pericyclic:electrocyclic}
$4n$ $\pi$ इलेक्ट्रॉनों वाले पॉलिईन की तापीय \omterm{def:b3:pericyclic:families}{विद्युतचक्रीय अभिक्रिया}
\omterm{def:b3:pericyclic:rotation-modes}{सहघूर्णी} होती है, और $4n + 2$ $\pi$ इलेक्ट्रॉनों वाले की \omterm{def:b3:pericyclic:rotation-modes}{विघूर्णी}।
\end{theorem}

\begin{proof}
नया $\sigma$ आबंध HOMO के अंतस्थ p कक्षकों से बनता है, जो
तापीय अभिक्रिया में उच्चतम ऊर्जा वाले इलेक्ट्रॉन रखता है; यह तभी आबंधी होता है जब
एक-दूसरे की ओर घूमने वाली पालियों का चिह्न समान हो। ह्युकेल सूत्र से
(\cref{ch:b3:solid-state}) कक्षक $k$ के गुणांक, $N$-परमाणु शृंखला के,
$\sin(jk\pi/(N + 1))$ के समानुपाती हैं, अतः दोनों सिरों, $j = 1$ और $j = N$, के
गुणांक $\sin(k\pi/(N + 1))$ और $\sin(Nk\pi/(N + 1)) = (-1)^{k+1}\sin(k\pi/(N + 1))$ हैं: विषम $k$ के लिए समान चिह्न,
सम $k$ के लिए विपरीत चिह्न। $N$ परमाणुओं और $N$ इलेक्ट्रॉनों के साथ HOMO
$k = N/2$ है। $N = 4n$ के लिए $k = 2n$ सम है: सिरों की ऊपरी पालियों के चिह्न
विपरीत हैं, और दोनों सिरों को एक ही दिशा में घूमना होगा, प्रत्येक अपनी
मेल खाते चिह्न वाली पालि को भीतर लाते हुए: \omterm{def:b3:pericyclic:rotation-modes}{सहघूर्णी}। $N = 4n + 2$ के लिए $k = 2n + 1$ विषम है: ऊपरी पालियों
का चिह्न समान है, और उन्हें एक-दूसरे की ओर घुमाना \omterm{def:b3:pericyclic:rotation-modes}{विघूर्णी} है।
\end{proof}

\begin{omfigure}
\begin{tikzpicture}[font=\small, >=Stealth]
  % butadiene HOMO psi2: 0.60 0.37 -0.37 -0.60 (lobes scaled x1.6)
  \foreach \x/\c in {0/0.60, 1/0.37, 2/-0.37, 3/-0.60} {
    \pic at (\x,0) {omp={1.6*\c}{90}};
    \fill (\x,0) circle (1.2pt);
  }
  \draw[black!60] (0,0) -- (3,0);
  \draw[->, thick, omThm] (-0.45,0.55) arc[start angle=130, end angle=20, radius=0.7];
  \draw[->, thick, omThm] (2.55,0.55) arc[start angle=130, end angle=20, radius=0.7];
  \node[font=\scriptsize, align=center] at (1.5,-1.3) {ब्यूटाडाईन, $\psi_2$ (HOMO, 4 इलेक्ट्रॉन):\\ विपरीत चिह्न के सिरे: \textbf{सहघूर्णी}};
  \begin{scope}[xshift=6.2cm]
  % hexatriene HOMO psi3: 0.52 0.23 -0.42 -0.42 0.23 0.52
  \foreach \x/\c in {0/0.52, 0.8/0.23, 1.6/-0.42, 2.4/-0.42, 3.2/0.23, 4.0/0.52} {
    \pic at (\x,0) {omp={1.6*\c}{90}};
    \fill (\x,0) circle (1.2pt);
  }
  \draw[black!60] (0,0) -- (4.0,0);
  \draw[->, thick, omThm] (-0.45,0.55) arc[start angle=130, end angle=20, radius=0.7];
  \draw[->, thick, omThm] (4.45,0.55) arc[start angle=50, end angle=160, radius=0.7];
  \node[font=\scriptsize, align=center] at (2.0,-1.3) {हेक्साट्राईन, $\psi_3$ (HOMO, 6 इलेक्ट्रॉन):\\ समान चिह्न के सिरे: \textbf{विघूर्णी}};
  \end{scope}
\end{tikzpicture}

{\small ब्यूटाडाईन और हेक्साट्राईन के HOMO, उनकी पालियाँ
ह्युकेल गुणांकों के अनुसार मापित (नीला और नारंगी: दोनों चिह्न), शृंखला के अनुदिश बनाई गईं।
सिरों के बीच आबंध बनने के लिए एक-दूसरे की ओर घूमने वाली पालियों का
चिह्न समान होना चाहिए: ब्यूटाडाईन में दोनों सिरे दक्षिणावर्त घूमते हैं, हेक्साट्राईन में विपरीत
दिशाओं में (लाल तीर)।}
\end{omfigure}

यह नियम त्रिविमविशिष्ट है। (\textit{E},\textit{E})-हेक्सा-2,4-डाईन, गर्म करने पर,
\omterm{def:b3:pericyclic:rotation-modes}{सहघूर्णी} रूप से \textit{trans}-3,4-डाइमेथिलसाइक्लोब्यूटीन में बंद होता; व्यवहार में तनावग्रस्त
साइक्लोब्यूटीन ही खुलता है, और \textit{cis}-3,4-डाइमेथिलसाइक्लोब्यूटीन
\omterm{def:b3:pericyclic:rotation-modes}{सहघूर्णी} रूप से केवल (\textit{E},\textit{Z})-हेक्सा-2,4-डाईन में खुलता है। (\textit{E},\textit{Z},\textit{E})-ऑक्टा-2,4,6-ट्राईन
\omterm{def:b3:pericyclic:rotation-modes}{विघूर्णी} रूप से \textit{cis}-5,6-डाइमेथिलसाइक्लोहेक्सा-1,3-डाईन में बंद होता है।

\begin{proposition}[प्रकाशरासायनिक विद्युतचक्रीय अभिक्रियाएँ]\label{prop:b3:pericyclic:photo-electrocyclic}
प्रकाश में नियम उलट जाते हैं: $4n$ इलेक्ट्रॉन, \omterm{def:b3:pericyclic:rotation-modes}{विघूर्णी}; $4n + 2$, \omterm{def:b3:pericyclic:rotation-modes}{सहघूर्णी}।
\end{proposition}

\begin{proof}
फ़ोटॉन का अवशोषण एक इलेक्ट्रॉन को HOMO $k = N/2$ से LUMO $k = N/2 + 1$ में उन्नत करता है
(\cref{ch:b3:photochemistry}), जो अब उच्चतम अधिकृत कक्षक है और
सिरों को नियंत्रित करता है। उसके $k$ की समता विपरीत है, अतः उसके
अंतस्थ गुणांकों का आपेक्षिक चिह्न उलट जाता है, और उसके साथ घूर्णन की दिशा भी।
\end{proof}

\begin{method}[विद्युतचक्रीय अभिक्रिया के त्रिविम रसायन का पूर्वानुमान]\label{met:b3:pericyclic:predict-stereo}
\begin{enumerate}
  \item विवृत-शृंखला पॉलिईन के $\pi$ इलेक्ट्रॉन गिनिए ($4n$ या $4n + 2$); ऊष्मा या
        प्रकाश नोट कीजिए।
  \item विधा पढ़िए: तापीय $4n$ सह, $4n + 2$ वि; प्रकाशरासायनिक उलटा।
  \item सिरों को उनके प्रतिस्थापियों सहित, बंद होने वाले संरूपण (s-सिस) में
        बनाइए; उन्हें चुनी गई विधा में 90° घुमाइए।
  \item पढ़िए कि दोनों प्रतिस्थापी वलय के एक ही फलक पर
        (\textit{cis}) समाप्त होते हैं या विपरीत फलकों पर (\textit{trans})। वलय-विवरण के लिए वही
        विश्लेषण उलटा चलाइए।
\end{enumerate}
\end{method}

\section{कक्षक सहसंबंध आरेख}

वुडवर्ड और हॉफ़मैन का मूल तर्क केवल HOMO को नहीं, सभी कक्षकों को
देखता है। \omterm{def:b3:pericyclic:rotation-modes}{सहघूर्णी} पथ के अनुदिश अणु एक $C_2$ अक्ष बनाए रखता है;
\omterm{def:b3:pericyclic:rotation-modes}{विघूर्णी} पथ के अनुदिश, एक \omterm{def:b3:point-groups:mirror}{दर्पण तल} $\sigma$। अभिकारक और
उत्पाद का प्रत्येक कक्षक बनाए रखे गए तत्व के सापेक्ष सममित (S) या प्रतिसममित (A)
होता है।

\begin{definition}[कक्षक सहसंबंध आरेख]\label{def:b3:pericyclic:orbital-correlation}
\emph{कक्षक सहसंबंध आरेख}\index{कक्षक सहसंबंध आरेख} अभिकारक के
प्रत्येक कक्षक को उत्पाद के उसी सममिति वाले कक्षक से जोड़ता है,
अभिक्रिया पथ के अनुदिश बनाए रखे गए \omterm{def:b3:point-groups:operation}{सममिति तत्वों} के सापेक्ष, उन्हें
ऊर्जा के क्रम में लेते हुए: सबसे निचला S सबसे निचले S से, इत्यादि।
\end{definition}

\begin{proposition}[कक्षक सममिति का संरक्षण]\label{prop:b3:pericyclic:correlation}
यदि अभिकारक में अधिकृत प्रत्येक कक्षक उत्पाद में अधिकृत किसी कक्षक से
सहसंबंधित हो, तो तापीय अभिक्रिया अनुमत है; यदि कोई अधिकृत कक्षक
रिक्त कक्षक से सहसंबंधित हो, तो वह निषिद्ध है, ऊँचे अवरोध के साथ।
\end{proposition}

\begin{proof}[तर्क]
पथ के अनुदिश कोई कक्षक अपनी सममिति बनाए रखता है और उसकी ऊर्जा सतत रूप से
बदलती है; समान सममिति के कक्षक एक-दूसरे को काट नहीं सकते (वे मिश्रित होकर प्रतिकर्षित होते हैं),
भिन्न सममिति के काट सकते हैं। जब अभिकारक के किसी द्वि-अधिकृत कक्षक को
उत्पाद का प्रतिआबंधी कक्षक बनना पड़ता है, तो अभिकारक की मूल अवस्था
उत्पाद की एक द्वि-उत्तेजित अवस्था तक ले जाती है, और मूल-अवस्था
ऊर्जा तेज़ी से तब तक बढ़ती है जब तक, पथ में बाद में, समान सममिति की अवस्थाएँ
मिश्रित न हों: अवरोध ऊँचा होता है।
\end{proof}

\begin{omfigure}
\begin{tikzpicture}[font=\scriptsize]
  \foreach \px/\title/\el in {0/विघूर्णी: $\sigma$ संरक्षित/s, 7.6/सहघूर्णी: $C_2$ संरक्षित/c} {
    \begin{scope}[xshift=\px cm]
    \node[font=\small\bfseries] at (2.0,5.1) {\title};
    \node at (0,4.7) {ब्यूटाडाईन};
    \node at (4.0,4.7) {साइक्लोब्यूटीन};
    \end{scope}
  }
  % ---- disrotatory (mirror): butadiene psi1 S, psi2 A, psi3 S, psi4 A;
  %      cyclobutene sigma S, pi S, pi* A, sigma* A
  \pic (b1) at (0,0.6) {omlevel={ud}{1.0}}; \node[anchor=east] at (-0.6,0.6) {$\psi_1$ S};
  \pic (b2) at (0,1.8) {omlevel={ud}{1.0}}; \node[anchor=east] at (-0.6,1.8) {$\psi_2$ A};
  \pic (b3) at (0,3.0) {omlevel={}{1.0}}; \node[anchor=east] at (-0.6,3.0) {$\psi_3$ S};
  \pic (b4) at (0,4.1) {omlevel={}{1.0}}; \node[anchor=east] at (-0.6,4.1) {$\psi_4$ A};
  \pic (c1) at (4.0,0.1) {omlevel={ud}{1.0}}; \node[anchor=west] at (4.6,0.1) {$\sigma$ S};
  \pic (c2) at (4.0,1.3) {omlevel={ud}{1.0}}; \node[anchor=west] at (4.6,1.3) {$\pi$ S};
  \pic (c3) at (4.0,3.4) {omlevel={}{1.0}}; \node[anchor=west] at (4.6,3.4) {$\pi^*$ A};
  \pic (c4) at (4.0,4.4) {omlevel={}{1.0}}; \node[anchor=west] at (4.6,4.4) {$\sigma^*$ A};
  \draw[omcorr] (b1-r) -- (c1-l);
  \draw[omcorrforbidden] (b2-r) -- (c3-l);
  \draw[omcorr] (b3-r) -- (c2-l);
  \draw[omcorr] (b4-r) -- (c4-l);
  % ---- conrotatory (C2): butadiene psi1 A, psi2 S, psi3 A, psi4 S;
  %      cyclobutene sigma S, pi A, pi* S, sigma* A
  \begin{scope}[xshift=7.6cm]
  \pic (d1) at (0,0.6) {omlevel={ud}{1.0}}; \node[anchor=east] at (-0.6,0.6) {$\psi_1$ A};
  \pic (d2) at (0,1.8) {omlevel={ud}{1.0}}; \node[anchor=east] at (-0.6,1.8) {$\psi_2$ S};
  \pic (d3) at (0,3.0) {omlevel={}{1.0}}; \node[anchor=east] at (-0.6,3.0) {$\psi_3$ A};
  \pic (d4) at (0,4.1) {omlevel={}{1.0}}; \node[anchor=east] at (-0.6,4.1) {$\psi_4$ S};
  \pic (e1) at (4.0,0.1) {omlevel={ud}{1.0}}; \node[anchor=west] at (4.6,0.1) {$\sigma$ S};
  \pic (e2) at (4.0,1.3) {omlevel={ud}{1.0}}; \node[anchor=west] at (4.6,1.3) {$\pi$ A};
  \pic (e3) at (4.0,3.4) {omlevel={}{1.0}}; \node[anchor=west] at (4.6,3.4) {$\pi^*$ S};
  \pic (e4) at (4.0,4.4) {omlevel={}{1.0}}; \node[anchor=west] at (4.6,4.4) {$\sigma^*$ A};
  \draw[omcorr] (d1-r) -- (e2-l);
  \draw[omcorr] (d2-r) -- (e1-l);
  \draw[omcorr] (d3-r) -- (e4-l);
  \draw[omcorr] (d4-r) -- (e3-l);
  \end{scope}
\end{tikzpicture}

{\small ब्यूटाडाईन $\rightleftharpoons$ साइक्लोब्यूटीन के लिए \omterm{def:b3:pericyclic:orbital-correlation}{कक्षक सहसंबंध आरेख}। बाएँ,
\omterm{def:b3:pericyclic:rotation-modes}{विघूर्णी} पथ एक \omterm{def:b3:point-groups:mirror}{दर्पण तल} बनाए रखता है: अधिकृत $\psi_2$ (A) साइक्लोब्यूटीन के
रिक्त $\pi^*$ (A) से सहसंबंधित होता है (लाल धराशायी): तापीय रूप से निषिद्ध। दाएँ,
\omterm{def:b3:pericyclic:rotation-modes}{सहघूर्णी} पथ एक $C_2$ अक्ष बनाए रखता है: अधिकृत कक्षक अधिकृत कक्षकों तक जाते हैं:
तापीय रूप से अनुमत।}
\end{omfigure}

\begin{method}[सहसंबंध आरेख बनाना]\label{met:b3:pericyclic:correlation-diagram}
\begin{enumerate}
  \item वे \omterm{def:b3:point-groups:operation}{सममिति तत्व} ज्ञात कीजिए जो बनने और टूटने वाले आबंधों को समद्विभाजित करते हैं और
        चुने गए पूरे पथ में बने रहते हैं।
  \item भाग लेने वाले अभिकारक और उत्पाद के कक्षकों को ऊर्जा के
        क्रम में सूचीबद्ध कीजिए, और प्रत्येक तत्व के लिए प्रत्येक को S या A नामांकित कीजिए।
  \item एक ओर के प्रत्येक सममिति प्रकार के सबसे निचले कक्षक को दूसरी ओर के उसी प्रकार के
        सबसे निचले से जोड़िए, फिर अगले को, समान सममिति की दो रेखाओं को कभी
        काटे बिना।
  \item अधिकृत कक्षकों की जाँच कीजिए: सभी अधिकृत तक, अनुमत; कोई एक रिक्त तक,
        निषिद्ध।
\end{enumerate}
\end{method}

\begin{definition}[अवस्था सहसंबंध आरेख]\label{def:b3:pericyclic:state-correlation}
\emph{अवस्था सहसंबंध आरेख}\index{अवस्था सहसंबंध आरेख} अभिकारक और उत्पाद की
इलेक्ट्रॉनिक अवस्थाओं को, जिनमें से प्रत्येक कक्षक विन्यासों से बनी और
अपनी समग्र सममिति से नामांकित होती है, पथ के माध्यम से जोड़ता है; समान सममिति की
अवस्थाएँ एक-दूसरे को नहीं काटतीं।
\end{definition}

निषिद्ध \omterm{def:b3:pericyclic:rotation-modes}{विघूर्णी} पथ के लिए ब्यूटाडाईन की मूल अवस्था ($\psi_1^2\psi_2^2$)
साइक्लोब्यूटीन की एक द्वि-उत्तेजित अवस्था ($\sigma^2\pi^{*2}$) से सहसंबंधित होती है, परंतु प्रथम उत्तेजित
अवस्था ($\psi_1^2\psi_2\psi_3$) साइक्लोब्यूटीन की प्रथम उत्तेजित अवस्था ($\sigma^2\pi\pi^*$) से सहसंबंधित होती है:
उत्तेजित अवस्था से \omterm{def:b3:pericyclic:rotation-modes}{विघूर्णी} पथ ढलान की ओर है। यही
\cref{prop:b3:pericyclic:photo-electrocyclic} का प्रकाशरासायनिक उलटाव है, अवस्थाओं पर देखा गया।

\section{चक्रीयोजन}

\begin{theorem}[चक्रीयोजन]\label{thm:b3:pericyclic:cycloaddition}
तापीय $[m + n]$ \omterm{def:b3:pericyclic:families}{चक्रीयोजन}, जिसमें दोनों घटक \omterm{def:b3:pericyclic:faciality}{समफलकीय} रूप से अभिक्रिया करते हैं,
तब अनुमत है जब $m + n = 4q + 2$ इलेक्ट्रॉन हों, और तब निषिद्ध जब $m + n = 4q$; प्रकाश में
नियम उलट जाता है।
\end{theorem}

\begin{proof}
दोनों नए आबंध एक घटक के HOMO और दूसरे के LUMO के बीच बनते हैं,
जो अपनी अंतस्थ पालियों से एक-दूसरे के आमने-सामने होते हैं; दोनों के
\omterm{def:b3:pericyclic:faciality}{समफलकीय} अभिक्रिया करने पर, दोनों आबंध तभी आबंधी होते हैं जब
HOMO और LUMO के अंतस्थ गुणांकों का आपेक्षिक चिह्न समान हो। $N$ परमाणुओं पर $N$
इलेक्ट्रॉनों वाले घटक के लिए HOMO का $k = N/2$ है और LUMO का $k = N/2 + 1$; अंतस्थ गुणांक
विषम $k$ के लिए समान चिह्न के, सम $k$ के लिए विपरीत। [4+2] स्थिति में डाईन HOMO ($k = 2$)
और एथीन LUMO ($k = 2$) दोनों के अंतस्थ चिह्न विपरीत हैं: वे मेल खाते हैं।
[2+2] स्थिति में एक एथीन का HOMO ($k = 1$, समान चिह्न) दूसरे के LUMO
($k = 2$, विपरीत) से मिलता है: एक सिरा आबंधी, एक प्रतिआबंधी। सामान्यतः $m$
घटक का HOMO और $n$ घटक का LUMO तब मेल खाते हैं जब $m/2$ और $n/2 + 1$ की
समता समान हो, अर्थात् जब $(m + n)/2$ विषम हो: $m + n = 4q + 2$। प्रकाश में उत्तेजित घटक का
एकल-अधिकृत LUMO उसके HOMO की भूमिका निभाता है, और समता शर्त
उलट जाती है।
\end{proof}

\begin{omfigure}
\begin{tikzpicture}[font=\scriptsize]
  % [4+2]: diene HOMO above, ethene LUMO (pi*) below, facing lobes match at both ends
  \foreach \x/\c in {0/0.60, 1/0.37, 2/-0.37, 3/-0.60} {
    \pic at (\x,1.6) {omp={1.4*\c}{90}}; \fill (\x,1.6) circle (1.1pt);
  }
  \draw[black!60] (0,1.6) -- (3,1.6);
  \foreach \x/\c in {0/-0.71, 3/0.71} {
    \pic at (\x,-0.4) {omp={1.4*\c}{90}}; \fill (\x,-0.4) circle (1.1pt);
  }
  \draw[black!60] (0,-0.4) -- (3,-0.4);
  \draw[omMeth, densely dashed, thick] (-0.25,0.75) -- (-0.25,0.35) (3.25,0.75) -- (3.25,0.35);
  \node[omMeth] at (-0.25,0.55) [left] {आबंधी};
  \node[omMeth] at (3.25,0.55) [right] {आबंधी};
  \node[font=\small] at (1.5,2.95) {[4+2]: डाईन HOMO};
  \node[font=\small] at (1.5,-1.75) {एथीन LUMO: अनुमत};
  \begin{scope}[xshift=7.2cm]
  \foreach \x/\c in {0/0.71, 3/0.71} {
    \pic at (\x,1.6) {omp={1.4*\c}{90}}; \fill (\x,1.6) circle (1.1pt);
  }
  \draw[black!60] (0,1.6) -- (3,1.6);
  \foreach \x/\c in {0/-0.71, 3/0.71} {
    \pic at (\x,-0.4) {omp={1.4*\c}{90}}; \fill (\x,-0.4) circle (1.1pt);
  }
  \draw[black!60] (0,-0.4) -- (3,-0.4);
  \draw[omMeth, densely dashed, thick] (-0.25,0.75) -- (-0.25,0.35);
  \draw[omThm, densely dashed, thick] (3.25,0.75) -- (3.25,0.35);
  \node[omMeth] at (-0.25,0.55) [left] {आबंधी};
  \node[omThm] at (3.25,0.55) [right] {प्रतिआबंधी};
  \node[font=\small] at (1.5,2.95) {[2+2]: एथीन HOMO};
  \node[font=\small] at (1.5,-1.75) {एथीन LUMO: निषिद्ध};
  \end{scope}
\end{tikzpicture}

{\small \omterm{def:b3:pericyclic:faciality}{समफलकीय} \omterm{def:b3:pericyclic:families}{चक्रीयोजनों} में आमने-सामने स्थित अग्रांत कक्षक (ऊपरी घटक की
निचली पालियाँ निचले घटक की ऊपरी पालियों से मिलती हैं; पालियों के
आकार ह्युकेल गुणांकों से)। [4+2] अभिक्रिया में दोनों नए आबंध
आबंधी हैं; [2+2] अभिक्रिया में एक सिरा आबंधी है और दूसरा प्रतिआबंधी।}
\end{omfigure}

तापीय रूप से दो ऐल्कीन एक समन्वित चरण में साइक्लोब्यूटेन नहीं देते;
प्रकाश में देते हैं (\cref{ch:b3:photochemistry})। कीटीन वह अपवाद हैं जो
नियम की पुष्टि करते हैं: उनका C=O $\pi^*$ कक्षक, जो C=C वाले के लंबवत है, उन्हें
किसी ऐल्कीन से [$\pi2_s$ + $\pi2_a$] ज्यामिति में अभिक्रिया करने देता है, एक घटक \omterm{def:b3:pericyclic:faciality}{समफलकीय} और
कीटीन \omterm{def:b3:pericyclic:faciality}{विपरीतफलकीय}, और वे तापीय रूप से साइक्लोब्यूटेनोन बनाते हैं।

\begin{definition}[1,3-द्विध्रुवी चक्रीयोजन]\label{def:b3:pericyclic:dipolar}
\emph{1,3-द्विध्रुव}\index{1,3-द्विध्रुव} चार इलेक्ट्रॉनों वाला एक त्रि-परमाणु $\pi$ निकाय है,
जो अपने सिरों पर या मध्य में औपचारिक आवेशों के साथ लिखा जाता है, जैसे
ऐज़ाइड $\ce{R-N=N+=N-}$, एक नाइट्राइल ऑक्साइड या एक ओज़ोन अणु। \emph{1,3-द्विध्रुवी
चक्रीयोजन}\index{1,3-द्विध्रुवी चक्रीयोजन} किसी ऐल्कीन या ऐल्काइन के साथ इसका [3+2] \omterm{def:b3:pericyclic:families}{चक्रीयोजन} है, एक
छह-इलेक्ट्रॉन अभिक्रिया, जो एक पंचसदस्यीय वलय देती है।
\end{definition}

किसी ऐज़ाइड की अंतस्थ ऐल्काइन के साथ कॉपर-उत्प्रेरित अभिक्रिया, जो 1,2,3-ट्राइऐज़ोल का
एक ही स्थिति-समावयवी देती है, क्लिक रसायन की विश्वसनीय जोड़ने वाली अभिक्रिया है।
ओज़ोनी-अपघटन ऐल्कीन पर ओज़ोन के \omterm{def:b3:pericyclic:dipolar}{1,3-द्विध्रुवी चक्रीयोजन} से आरंभ होता है।

\section{सिग्माट्रॉपी पुनर्विन्यास और व्यापक नियम}

\begin{proposition}[हाइड्रोजन विस्थापन]\label{prop:b3:pericyclic:sigmatropic-h}
किसी पॉलिईन के अनुदिश हाइड्रोजन का तापीय $[1, j]$ विस्थापन \omterm{def:b3:pericyclic:faciality}{समफलकीय} रूप से अनुमत है जब
संक्रमण अवस्था में $4q + 2$ इलेक्ट्रॉन हों ($[1,5]$), और \omterm{def:b3:pericyclic:faciality}{विपरीतफलकीय} होना आवश्यक है जब उसमें $4q$ हों
($[1,3]$, $[1,7]$)।
\end{proposition}

\begin{proof}
संक्रमण अवस्था को ऐसे हाइड्रोजन परमाणु (1s) के रूप में लीजिए जो $j$ परमाणुओं और
$j$ इलेक्ट्रॉनों वाले मूलक के सिरों के बीच गति करता है; हाइड्रोजन SOMO, $k = (j + 1)/2$, के दोनों सिरों को
जोड़ता है। $j = 5$ के लिए $k = 3$ विषम है, सिरों की समान चिह्न वाली पालियाँ एक ही फलक पर हैं,
और 1s कक्षक उस फलक पर दोनों से आबंधित हो सकता है: \omterm{def:b3:pericyclic:faciality}{समफलकीय}। $j = 3$ और $j = 7$ के लिए
$k = 2$ और $4$ सम हैं, समान-चिह्न पालियाँ दोनों सिरों पर विपरीत फलकों पर हैं, और
हाइड्रोजन को एक फलक से दूसरे पर जाना पड़ता है: \omterm{def:b3:pericyclic:faciality}{विपरीतफलकीय}। $[1,3]$ विस्थापन
इतनी छोटी शृंखला के पार नहीं पहुँच सकता; $[1,7]$ विस्थापन एक कुंडलित पॉलिईन के अनुदिश पहुँच सकता है।
\end{proof}

\begin{definition}[कोप और क्लाइज़न पुनर्विन्यास]\label{def:b3:pericyclic:cope-claisen}
\emph{कोप पुनर्विन्यास}\index{कोप पुनर्विन्यास} हेक्सा-1,5-डाईन का [3,3] \omterm{def:b3:pericyclic:families}{सिग्माट्रॉपी पुनर्विन्यास} है;
\emph{क्लाइज़न
पुनर्विन्यास}\index{क्लाइज़न पुनर्विन्यास} किसी ऐलिल वाइनिल ईथर, या
किसी ऐलिल ऐरिल ईथर, का $\gamma,\delta$-असंतृप्त कार्बोनिल यौगिक या \textit{ऑर्थो}-ऐलिलफ़ीनॉल में ऐसा पुनर्विन्यास।
\end{definition}

\begin{omfigure}
\begin{tikzpicture}[font=\scriptsize, scale=1.1]
  \foreach \px/\three/\title in {0/C/कोप, 5.0/O/क्लाइज़न} {
    \begin{scope}[xshift=\px cm]
      % chair: atoms at 60-degree steps, alternately 0.25 above and below the
      % mean plane, seen from slightly above (y' = 0.35 y + z), scaled
      \begin{scope}[shift={(1.3,0.2)}, scale=1.15]
      \coordinate (a1) at (-1.25,-0.25); \coordinate (a2) at (-0.625,-0.129); \coordinate (a3) at (0.625,-0.629);
      \coordinate (a4) at (1.25,0.25); \coordinate (a5) at (0.625,0.129); \coordinate (a6) at (-0.625,0.629);
      \end{scope}
      \draw[thick] (a1) -- (a2) -- (a3);
      \draw[thick] (a4) -- (a5) -- (a6);
      \draw[omThm, thick, densely dashed] (a3) -- (a4);
      \draw[omMeth, thick, densely dashed] (a6) -- (a1);
      \foreach \i/\l in {1/1, 2/2, 4/4, 5/5, 6/6} { \fill (a\i) circle (1.3pt); }
      \node[fill=white, inner sep=0.6pt] at (a3) {\three};
      \node[anchor=east] at (a1) {1};
      \node[anchor=north] at (a2) {2};
      \node[anchor=north west] at (a3) {3};
      \node[anchor=west] at (a4) {4};
      \node[anchor=south] at (a5) {5};
      \node[anchor=south] at (a6) {6};
      \node[font=\small\bfseries] at (1.3,1.5) {\title};
    \end{scope}
  }
\end{tikzpicture}

{\small [3,3] \omterm{def:b3:pericyclic:families}{सिग्माट्रॉपी पुनर्विन्यासों} की कुर्सी संक्रमण अवस्थाएँ: दो ऐलिल
अणुखंड आमने-सामने हैं; 3--4 $\sigma$ आबंध टूटता है (लाल) जबकि 1--6 आबंध बनता है
(हरा)। \omterm{def:b3:pericyclic:cope-claisen}{क्लाइज़न पुनर्विन्यास} में परमाणु 3 ईथर ऑक्सीजन है, और
उत्पाद एक कार्बोनिल यौगिक है। कुर्सी नए त्रिविमजनक केंद्रों का आपेक्षिक विन्यास
नियत करती है।}
\end{omfigure}

[3,3] विस्थापन एक छह-इलेक्ट्रॉन, पूर्णतः \omterm{def:b3:pericyclic:faciality}{समफलकीय} प्रक्रम है, जो तापीय रूप से अनुमत है;
यह साइक्लोहेक्सेन की तरह कुर्सी-सदृश संक्रमण अवस्था से होकर चलता है, और
प्रतिस्थापी उसकी विषुवतीय स्थितियाँ पसंद करते हैं: (\textit{E},\textit{E}) और (\textit{Z},\textit{Z}) डाईन
पूर्वानुमेय रूप से विपरीत आपेक्षिक विन्यास देते हैं।

\begin{theorem}[वुडवर्ड--हॉफ़मैन नियम]\label{thm:b3:pericyclic:woodward-hoffmann}
कोई तापीय \omterm{def:b3:pericyclic:pericyclic}{परिचक्रीय अभिक्रिया} तब अनुमत है जब $(4q + 2)_s$ और
$(4r)_a$ घटकों की कुल संख्या विषम हो; प्रकाशरासायनिक तब, जब वह सम हो।
\end{theorem}

\begin{proof}[आंशिक उपपत्ति]
ऊपर विवेचित प्रत्येक परिवार के लिए गणना अग्रांत-कक्षक
विश्लेषण से सहमत है: तापीय [4+2] \omterm{def:b3:pericyclic:families}{चक्रीयोजन} के घटक $\pi4_s$ और $\pi2_s$ हैं, एक $(4q + 2)_s$
घटक, विषम: अनुमत; [2+2] में $\pi2_s + \pi2_s$ है, दो, सम: निषिद्ध, जबकि
$\pi2_s + \pi2_a$ में एक: अनुमत; साइक्लोब्यूटीन का \omterm{def:b3:pericyclic:rotation-modes}{सहघूर्णी} विवरण $\sigma2_s + \pi2_a$ है, एक:
अनुमत; \omterm{def:b3:pericyclic:faciality}{समफलकीय} $[1,5]$-H विस्थापन $\sigma2_s + \pi4_s$ है, जिसमें केवल $\sigma2_s$ गिना जाता है: एक,
अनुमत। व्यापक उपपत्ति, जो घटकों के हर विन्यास के लिए मान्य है, अधिक उन्नत
पाठ्यक्रमों का विषय है।
\end{proof}

\begin{definition}[ऐरोमैटिक और मोबियस संक्रमण अवस्थाएँ]\label{def:b3:pericyclic:mobius}
\emph{ऐरोमैटिक संक्रमण अवस्था}\index{ऐरोमैटिक संक्रमण अवस्था} किसी परिचक्रीय संक्रमण
अवस्था में कक्षकों का ऐसा चक्रीय विन्यास है जो, ह्युकेल वलय की तरह,
$4q + 2$ इलेक्ट्रॉन धारण करने पर स्थायी होता है। \emph{मोबियस संक्रमण
अवस्था}\index{मोबियस संक्रमण अवस्था} में वलय के चारों ओर पड़ोसी कक्षकों के बीच
चिह्न परिवर्तनों की संख्या विषम होती है (एक \omterm{def:b3:pericyclic:faciality}{विपरीतफलकीय} ऐंठन); यह
$4q$ इलेक्ट्रॉनों के साथ स्थायी होती है।
\end{definition}

\begin{proposition}[ऐरोमैटिक संक्रमण अवस्थाएँ]\label{prop:b3:pericyclic:aromatic-ts}
कोई तापीय \omterm{def:b3:pericyclic:pericyclic}{परिचक्रीय अभिक्रिया} तब अनुमत है जब उसकी संक्रमण अवस्था ऐरोमैटिक हो:
$4q + 2$ इलेक्ट्रॉनों के साथ ह्युकेल सांस्थिति, या $4q$ इलेक्ट्रॉनों के साथ मोबियस सांस्थिति।
\end{proposition}

\begin{proof}[तर्क]
संक्रमण अवस्था में कक्षकों का चक्रीय विन्यास किसी चक्रीय
संयुग्मित अणु की तरह व्यवहार करता है। ह्युकेल वलय का निम्नतम स्तर अनपभ्रष्ट होता है और उसके ऊपर
युग्म, जो $4q + 2$ इलेक्ट्रॉनों से भरते हैं (द्वितीय वर्ष का खंड); एक चिह्न
प्रतिलोमन वाले वलय के सभी स्तर युग्मों में होते हैं, जो $4q$ से भरते हैं। भरा कोश अर्थात्
स्थायीकरण, एक \omterm{def:b3:pericyclic:mobius}{ऐरोमैटिक संक्रमण अवस्था} और निम्न अवरोध। चिह्न
परिवर्तनों की गिनती \omterm{def:b3:pericyclic:faciality}{विपरीतफलकीय} घटकों की गिनती है, अतः यह नियम
वुडवर्ड--हॉफ़मैन नियम के तुल्य है।
\end{proof}

\begin{method}[घटक बनाना और गिनना]\label{met:b3:pericyclic:components}
\begin{enumerate}
  \item बनने और टूटने वाले आबंध चिह्नित कीजिए; बदलने वाला प्रत्येक $\pi$ निकाय या $\sigma$ आबंध एक
        घटक है, जिसे उसके प्रकार और इलेक्ट्रॉन गणना से नामांकित किया जाता है ($\pi4$, $\pi2$, $\sigma2$, रिक्त कक्षक के लिए
        $\omega0$)।
  \item ज्यामिति (एक ही फलक, या
        विपरीत फलक) से प्रत्येक घटक के लिए s या a चुनिए।
  \item $(4q + 2)_s$ और $(4r)_a$ घटक गिनिए: विषम, तापीय रूप से अनुमत; सम,
        प्रकाशरासायनिक रूप से अनुमत।
  \item वलय के चारों ओर इलेक्ट्रॉन गणना से प्रति-जाँच कीजिए: $4q + 2$ के साथ ह्युकेल,
        $4q$ के साथ मोबियस।
\end{enumerate}
\end{method}

\begin{omfigure}
\begin{tikzpicture}[font=\scriptsize, >=Stealth]
  \foreach \px/\mode in {0/1, 4.6/2, 9.2/3} {
    \begin{scope}[xshift=\px cm]
      \foreach \n/\a in {10/210, 5/150, 6/90, 7/30, 8/-30, 9/-90} { \coordinate (v\n) at (\a:0.85); }
      \ifnum\mode=1
        \draw[thick] (v10) -- (v5) (v6) -- (v7) (v8) -- (v9);
        \draw[thick, double, double distance=1.6pt] (v5) -- (v6) (v7) -- (v8);
        \draw[thick, omThm] (v9) -- (v10);
        \draw[thick] (v10) -- ++(240:0.6) node[anchor=north] {$\ce{CH3}$};
        \node[font=\small\bfseries] at (0,1.55) {7-डीहाइड्रोकोलेस्टेरॉल};
      \fi
      \ifnum\mode=2
        \draw[thick, double, double distance=1.6pt] (v10) -- (v5) (v6) -- (v7) (v8) -- (v9);
        \draw[thick] (v5) -- (v6) (v7) -- (v8);
        \draw[thick] (v10) -- ++(240:0.6) node[anchor=north] (me) {$\ce{CH3}$};
        \draw[->, omThm, thick] ($(v10)+(240:0.75)+(0.35,0)$) to[bend right=40] ($(v9)+(-0.12,-0.12)$);
        \node[font=\small\bfseries] at (0,1.55) {प्रीविटामिन D$_3$};
      \fi
      \ifnum\mode=3
        \draw[thick] (v10) -- (v5) (v6) -- (v7) (v8) -- (v9);
        \draw[thick, double, double distance=1.6pt] (v5) -- (v6) (v7) -- (v8);
        \draw[thick, double, double distance=1.6pt] (v10) -- ++(240:0.6) node[anchor=north] {$\ce{CH2}$};
        \node[anchor=north west, inner sep=1pt] at (v9) {$\ce{CH2}$};
        \node[font=\small\bfseries] at (0,1.55) {विटामिन D$_3$};
      \fi
      \foreach \n/\a in {10/210, 5/150, 6/90, 7/30, 8/-30, 9/-110} {
        \node[font=\tiny, black!60] at (\a:1.12) {\n};
      }
    \end{scope}
  }
  \draw[->, thick] (1.45,0) -- (3.0,0) node[midway, above, align=center] {$h\nu$\\ 6 e$^-$, सहघूर्णी};
  \draw[->, thick] (6.05,0) -- (7.6,0) node[midway, above, align=center] {ऊष्मा, $[1,7]$-H\\ विपरीतफलकीय};
\end{tikzpicture}

{\small विटामिन D संश्लेषण के दो परिचक्रीय चरण, अभिक्रिया करने वाले परमाणुओं पर
दिखाए गए (धूसर में स्टेरॉयड क्रमांकन; अणु का शेष भाग, वलय A, C और D,
छोड़ दिया गया)। प्रकाश 9--10 आबंध (लाल) तोड़कर वलय B खोलता है, एक प्रकाशरासायनिक
छह-इलेक्ट्रॉन \omterm{def:b3:pericyclic:families}{विद्युतचक्रीय अभिक्रिया}, \omterm{def:b3:pericyclic:rotation-modes}{सहघूर्णी}। प्रीविटामिन D$_3$ में फिर
C19 मेथिल समूह का एक हाइड्रोजन सात-परमाणु कुंडलित ट्राईन के पार C9 पर (लाल तीर)
\omterm{def:b3:pericyclic:faciality}{विपरीतफलकीय} रूप से खिसकता है।}
\end{omfigure}

\begin{inthelab}[एक क्लाइज़न पुनर्विन्यास]
एक ऐलिल ऐरिल ईथर को नाइट्रोजन के अधीन पश्चवाह पर, 1,2-डाइक्लोरोबेंज़ीन जैसे
उच्च-क्वथनांकी विलायक में, या बिना विलायक के, लगभग \qty{200}{\celsius} पर गर्म किया जाता है, और
अभिक्रिया का अनुसरण पतली-परत वर्णलेखन से किया जाता है: बना \textit{ऑर्थो}-ऐलिलफ़ीनॉल
ईथर से अधिक ध्रुवीय है। फिर फ़ीनॉल को जलीय सोडियम
हाइड्रॉक्साइड में निष्कर्षित करके उदासीन अशुद्धियों से पृथक किया जाता है, और अम्लीकरण द्वारा पुनः प्राप्त किया जाता है।
\end{inthelab}

\begin{safety}{\ghs[0.62cm]{GHS02}\ghs[0.62cm]{GHS04}\ghs[0.62cm]{GHS08}}
% ledger: ghs4:butadiene
ब्यूटा-1,3-डाईन, एक सामान्य डाईन: दाब में अत्यधिक ज्वलनशील गैस, आनुवंशिक दोष और
कैंसर उत्पन्न कर सकती है। इसे केवल बंद प्रणालियों में सँभाला जाता है; शिक्षण में
ठोस सल्फ़ोलीन, जो गर्म करने पर \omterm{def:b3:pericyclic:families}{कीलेट्रॉपी अभिक्रिया} में इसे मुक्त करता है,
इसके स्थान पर एक बंद उपकरण में प्रयुक्त होता है।
\end{safety}

\begin{history}[कक्षक सममिति]
1965 में रॉबर्ट वुडवर्ड, जिन्होंने विटामिन B$_{12}$ का संश्लेषण करते हुए विद्युतचक्रीय चरणों का
उलझन भरा त्रिविम रसायन पाया था, और रोआल्ड हॉफ़मैन ने कक्षक सममिति के संरक्षण के
नियम प्रकाशित किए। केनिची फ़ुकुई ने 1952 में अग्रांत कक्षक
प्रस्तुत किए थे; फ़ुकुई और हॉफ़मैन ने रसायन का 1981 का नोबेल पुरस्कार
साझा किया। वुडवर्ड, जिन्हें अपने संश्लेषणों के लिए 1965 का पुरस्कार मिला था, का
1979 में निधन हुआ।
\end{history}

\section{अभ्यास}

\begin{exercise}[$\star$]\label{exo:b3:pericyclic:1}
वर्गीकृत कीजिए: डील्स--ऐल्डर अभिक्रिया; साइक्लोब्यूटीन का वलय-विवरण; \omterm{def:b3:pericyclic:cope-claisen}{कोप
पुनर्विन्यास}; सल्फ़ोलीन से $\ce{SO2}$ का निकास; किसी ऐल्कीन से ओज़ोन का योग;
साइक्लोपेन्टाडाईन में $[1,5]$-H विस्थापन।
\end{exercise}

\begin{exercise}[$\star$]\label{exo:b3:pericyclic:2}
हेक्सा-1,3,5-ट्राईन और ब्यूटा-1,3-डाईन के तापीय और
प्रकाशरासायनिक वलय समापन की विधा, \omterm{def:b3:pericyclic:rotation-modes}{सहघूर्णी} या \omterm{def:b3:pericyclic:rotation-modes}{विघूर्णी}, का पूर्वानुमान कीजिए।
\end{exercise}

\begin{exercise}[$\star$]\label{exo:b3:pericyclic:3}
\textit{trans}-3,4-डाइमेथिलसाइक्लोब्यूटीन गर्म करने पर कौन-सा उत्पाद देता है?
\end{exercise}

\begin{exercise}[$\star$]\label{exo:b3:pericyclic:4}
एथीन गर्म करने पर साइक्लोब्यूटेन में द्वितयित क्यों नहीं होता, जबकि प्रकाश में होता है
(सुग्राहीकारक के साथ या प्रत्यक्ष उत्तेजन द्वारा)?
\end{exercise}

\begin{exercise}[$\star\star$]\label{exo:b3:pericyclic:5}
\omterm{def:b3:pericyclic:rotation-modes}{विघूर्णी} पथ के लिए हेक्साट्राईन $\rightleftharpoons$ साइक्लोहेक्सा-1,3-डाईन का \omterm{def:b3:pericyclic:orbital-correlation}{कक्षक सहसंबंध आरेख}
बनाइए, और निष्कर्ष निकालिए।
\end{exercise}

\begin{exercise}[$\star\star$]\label{exo:b3:pericyclic:6}
प्रकाश में (\textit{E},\textit{Z},\textit{E})-ऑक्टा-2,4,6-ट्राईन किस डाइमेथिलसाइक्लोहेक्साडाईन में बंद होता है?
\end{exercise}

\begin{exercise}[$\star\star$]\label{exo:b3:pericyclic:7}
5-मेथिलसाइक्लोपेन्टा-1,3-डाईन कक्ष ताप पर 1- और
2-मेथिलसाइक्लोपेन्टाडाईनों में पुनर्विन्यासित होता है। $[1,5]$-H विस्थापन से समझाइए, और बताइए कि यह \omterm{def:b3:pericyclic:faciality}{समफलकीय} क्यों है।
\end{exercise}

\begin{exercise}[$\star\star$]\label{exo:b3:pericyclic:8}
किसी ऐल्कीन पर कीटीन के \omterm{def:b3:pericyclic:families}{चक्रीयोजन} के घटक गिनिए, और दिखाइए
कि यह तापीय रूप से अनुमत है।
\end{exercise}

\begin{exercise}[$\star\star$]\label{exo:b3:pericyclic:9}
ऐलिल वाइनिल ईथर, $\ce{CH2=CH-O-CH2-CH=CH2}$, गर्म करने पर \omterm{def:b3:pericyclic:cope-claisen}{क्लाइज़न पुनर्विन्यास} से
गुज़रता है। इसकी कुर्सी संक्रमण अवस्था बनाइए और उत्पाद बताइए।
\end{exercise}

\begin{exercise}[$\star\star\star$]\label{exo:b3:pericyclic:10}
एक संक्रमण अवस्था में एक \omterm{def:b3:pericyclic:faciality}{विपरीतफलकीय} घटक के साथ चक्रीय विन्यास में आठ इलेक्ट्रॉन
हैं। क्या यह ऐरोमैटिक है? क्या अभिक्रिया तापीय रूप से अनुमत है?
\end{exercise}

\begin{exercise}[$\star\star\star$]\label{exo:b3:pericyclic:11}
किसी अंतस्थ ऐल्काइन के साथ किसी ऐज़ाइड के \omterm{def:b3:pericyclic:dipolar}{1,3-द्विध्रुवी चक्रीयोजन} में ऐज़ाइड का सबसे बड़ा
HOMO गुणांक अंतस्थ नाइट्रोजन N3 पर है और ऐल्काइन का सबसे बड़ा
LUMO गुणांक उसके अंतस्थ कार्बन पर (अभ्यास के आँकड़े)।
सबसे बड़े गुणांकों को युग्मित करके, उत्प्रेरक के बिना अनुकूल स्थिति-समावयवी का पूर्वानुमान कीजिए, और
टिप्पणी कीजिए।
\end{exercise}

\begin{exercise}[$\star\star\star$]\label{exo:b3:pericyclic:12}
ह्युकेल सूत्र का उपयोग करके दिखाइए कि $[1,7]$-H विस्थापन \omterm{def:b3:pericyclic:faciality}{विपरीतफलकीय} होना चाहिए,
और $[1,5]$-H विस्थापन \omterm{def:b3:pericyclic:faciality}{समफलकीय}।
\end{exercise}

% keep the heading with its problem box, which fills a page by itself
\clearpage
\enlargethispage{3\baselineskip}
\section{समस्या: सूर्यप्रकाश, त्वचा और विटामिन D}

\begin{problem}[{सप्ताहांत समस्या --- सूर्यप्रकाश, त्वचा और विटामिन D:
7-डीहाइड्रोकोलेस्टेरॉल का प्रकाशरासायनिक वलय-विवरण, तापीय हाइड्रोजन विस्थापन,
पार्श्व उत्पाद, और शरीर को यह कार्य पूरा करने में कितना समय लगता है}]
\label{pb:b3:pericyclic:1}
समस्या के आँकड़े: \qty{37}{\celsius} पर प्रीविटामिन D$_3$ का विटामिन D$_3$ में रूपांतरण
$k = \qty{2.3e-2}{h^{-1}}$ के साथ प्रथम कोटि का है (इसके उलटे की उपेक्षा की गई है); इसकी सक्रियण ऊर्जा
\qty{85}{kJ/mol} है।

\medskip\noindent\textbf{भाग I --- वलय-विवरण।}
\begin{enumerate}
  \item 7-डीहाइड्रोकोलेस्टेरॉल का कौन-सा आबंध टूटता है, और कौन-सा $\pi$ निकाय बनता है?
  \item कितने इलेक्ट्रॉन भाग लेते हैं, और अभिक्रिया का परिवार क्या है?
  \item इसे प्रकाश की आवश्यकता क्यों है?
  \item प्रकाशरासायनिक विवरण \omterm{def:b3:pericyclic:rotation-modes}{सहघूर्णी} है या \omterm{def:b3:pericyclic:rotation-modes}{विघूर्णी}? कक्षकों से
        औचित्य बताइए।
  \item प्रीविटामिन D$_3$ का नया केंद्रीय द्विबंध \textit{Z} क्यों होना चाहिए?
  \item शरीर के ताप पर अँधेरे में वही विवरण तापीय रूप से क्यों नहीं हो सकता?
\end{enumerate}

\medskip\noindent\textbf{भाग II --- हाइड्रोजन विस्थापन।}
\begin{enumerate}[resume]
  \item उन परमाणुओं के नाम बताइए जिनके बीच हाइड्रोजन खिसकता है, और विस्थापन का प्रकार।
  \item संक्रमण अवस्था में कितने इलेक्ट्रॉन भाग लेते हैं?
  \item घटक गिनिए और निर्णय कीजिए कि तापीय अभिक्रिया
        \omterm{def:b3:pericyclic:faciality}{समफलकीय} रूप से अनुमत है या \omterm{def:b3:pericyclic:faciality}{विपरीतफलकीय} रूप से।
  \item यहाँ \omterm{def:b3:pericyclic:faciality}{विपरीतफलकीय} विस्थापन ज्यामितीय रूप से संभव क्यों है, जबकि $[1,3]$ विस्थापन
        नहीं?
  \item संक्रमण अवस्था ह्युकेल है या मोबियस, और क्या वह ऐरोमैटिक है?
  \item विस्थापन को प्रकाश की आवश्यकता क्यों नहीं है?
\end{enumerate}

\medskip\noindent\textbf{भाग III --- पार्श्व उत्पाद।}
\begin{enumerate}[resume]
  \item प्रीविटामिन D$_3$ प्रकाश में फिर से बंद हो सकता है। किस विधा में, और यह
        7-डीहाइड्रोकोलेस्टेरॉल का एक वलय त्रिविम समावयवी (ल्यूमिस्टेरॉल) क्यों दे सकता है?
  \item प्रकाश केंद्रीय \textit{Z} द्विबंध को \textit{E} (टैकिस्टेरॉल) में भी बदल सकता है।
        टैकिस्टेरॉल $[1,7]$-H विस्थापन से क्यों नहीं गुज़र सकता?
  \item लंबा सूर्यप्रकाश अधिकाधिक विटामिन D क्यों नहीं बनाता?
  \item वलय-विवरण तीव्र पर हाइड्रोजन विस्थापन धीमा क्यों है?
\end{enumerate}

\medskip\noindent\textbf{भाग IV --- शरीर के ताप पर बलगतिकी।}
\begin{enumerate}[resume]
  \item \qty{37}{\celsius} पर प्रीविटामिन D$_3$ की अर्ध-आयु परिकलित कीजिए।
  \item एक दिन में रूपांतरित अंश परिकलित कीजिए।
  \item 90\,\% रूपांतरण का समय परिकलित कीजिए।
  \item \qty{20}{\celsius} पर $k$ परिकलित कीजिए।
  \item \qty{20}{\celsius} पर 90\,\% रूपांतरण का समय परिकलित कीजिए।
  \item त्वचा का गर्म होना क्यों उपयोगी है?
  \item विटामिन D पूरक किसी पारंपरिक संश्लेषण के बजाय किसी स्टेरॉल को विकिरित करके और गर्म करके
        क्यों बनाए जाते हैं?
  \item परिणाम बताइए: शरीर के ताप पर प्रीविटामिन D$_3$ के 90\,\% रूपांतरण का समय।
\end{enumerate}
\end{problem}
